\chapter{اللقاء الرابع والثلاثون: الطلاق قبل المسيس والمتعة ونصف الصداق}

\section{مقدمة}
السلام عليكم ورحمة الله وبركاته. الحمد لله رب العالمين، وأشهد أن لا إله إلا الله وحده لا شريك له، وأشهد أن محمداً عبده ورسوله، صلى الله عليه وعلى آله وصحبه وسلم.

استأنف الشيخ هذا اللقاء من قوله تعالى: \begin{ayah}لَا جُنَاحَ عَلَيْكُمْ إِنْ طَلَّقْتُمُ النِّسَاءَ مَا لَمْ تَمَسُّوهُنَّ أَوْ تَفْرِضُوا لَهُنَّ فَرِيضَةً\end{ayah}، فتكلم على أحكام الطلاق قبل الدخول، والمتعة، ثم دخل في قوله تعالى: \begin{ayah}وَإِنْ طَلَّقْتُمُوهُنَّ مِنْ قَبْلِ أَنْ تَمَسُّوهُنَّ وَقَدْ فَرَضْتُمْ لَهُنَّ فَرِيضَةً\end{ayah}، وما فيها من أحكام نصف الصداق والعفو، ومن هو المراد بقوله: \begin{ayah}أَوْ يَعْفُوَا الَّذِي بِيَدِهِ عُقْدَةُ النِّكَاحِ\end{ayah}.

وهذا المجلس من أنفع المجالس في إحكام التصوير؛ لأنه يبين أن اختلاف صور الطلاق قبل الدخول يوجب اختلاف الأحكام، وأن القرآن يفصل هذه الصور تفصيلاً بديعاً.

\separator

\section{تأويل (لا جناح عليكم إن طلقتم النساء ما لم تمسوهن)}
\isnad{قال أبو جعفر:} نفى الله الجناح عن طلاق النساء قبل المسيس في هذه الصورة، وفسر \rawi{الطبري} \textit{المس} هنا بالجماع، كما هو المعروف في لسان القرآن، مع استعمال الكناية الأدبِيَّة التي يكثر ورودها في هذا الباب.

ثم تعرض لاختلاف القراء في \textit{تمسوهن} و\textit{تماسوهن}، وقرر أن القراءتين صحيحتان متفقتا المعنى في الجملة، وإن كان في إحداهما إشعار بالفعل من الجانبين على وجه أظهر. لكنه نبّه إلى أن هذه الزيادة في المعنى لا توجب اختلافاً في الحكم، فبقيت القراءتان جميعاً على الصحة والاعتبار.

\begin{fawaid}[ليس كل زيادة في المعنى توجب اختلافاً في الحكم]
قد تختلف القراءتان في جهة من جهات الدلالة، لكن يبقى الحكم الشرعي المستفاد منهما واحداً. ومن فقه الطبري أنه لا يرد القراءة الصحيحة لمجرد وجود زيادة معنى لا يترتب عليها اختلاف الحكم.
\end{fawaid}

\section{تحرير صورة المرأة المعنية بالآية}
قرر \rawi{الطبري} أن الآية تتناول في أصلها صورتين:
\begin{itemize}
    \item امرأة طلقت قبل المسيس ولم يسم لها صداق.
    \item وامرأة طلقت قبل المسيس مع فرض الصداق، ويأتي حكمها مفصلاً بعد ذلك.
\end{itemize}

وهذا من دقة القرآن في التفصيل؛ إذ لم يجمع الصورتين في حكم واحد، بل فصل بينهما، لأن ما وجب في ذمة الزوج في كل منهما مختلف. ومن هنا تظهر عناية \rawi{الطبري} بتتبع نظم الآية، وبيان أن ما عطف بعضه على بعض ليس تكراراً محضاً، بل بيان لصورتين وحكمين.

\begin{fawaid}[القرآن يفرق بين الصور المتقاربة إذا اختلف أثرها]
قد تبدو المسألتان متقاربتين في الظاهر، لكن اختلاف وجود الصداق أو عدمه يوجب اختلاف الحكم، ولذلك جاء البيان مفصلاً. ومن أحسن الفقه أن لا يسوى بين صورتين فرق الشارع بينهما.
\end{fawaid}

\section{تأويل (أو تفرضوا لهن فريضة)}
فسر \rawi{الطبري} الفريضة هنا بالصداق الواجب، وبيّن أن أصل الفرض في اللغة هو الإيجاب والإلزام. فمعنى الآية: أو توجبوا لهن صداقاً مسمى معلوماً.

ومن هذا الموضع يظهر كيف يربط \rawi{الطبري} بين أصل اللفظ في العربية وبين الحكم الفقهي المستفاد منه، فلا يفصل أحد الجانبين عن الآخر.

\separator

\section{تأويل (ومتعوهن على الموسع قدره وعلى المقتر قدره)}
انتقل \rawi{الطبري} بعد ذلك إلى المتعة التي تستحقها المطلقة قبل المسيس إذا لم يفرض لها صداق، فذكر أقوال السلف في مقدارها:
\begin{itemize}
    \item فقيل: أعلاها الخادم، ثم الورق، ثم الكسوة.
    \item وقيل: يراعى فيها الوسط والعدل بحسب حال الزوج.
    \item وقيل: يرجع فيها إلى اجتهاد الحاكم عند التنازع.
\end{itemize}

ثم ظهر من مجموع صنيعه أن أصل الحكم فيها أنها متعة معتبرة بحسب سعة الزوج وإقتاره، لا بمقدار جامد يلزم كل الناس على السواء. فالآية نفسها ردت الأمر إلى الموسع والمقتر، وجعلت المعروف هو الضابط.

\begin{fawaid}[المقادير التي يختلف فيها الناس يرد كثير منها إلى المعروف]
ليس كل حق مالي في الشريعة مقدراً تقديراً واحداً لا يختلف. فمنها ما يرد إلى الأعراف العادلة وأحوال الناس، كما في المتعة والنفقة ونحو ذلك، ما دام الضابط هو المعروف والعدل.
\end{fawaid}

\section{معنى (لا جناح عليكم) في هذا الموضع}
ذكر الشيخ أن \rawi{الطبري} لم يحمل نفي الجناح هنا على مجرد الإباحة العامة فحسب، بل وجّه بعض الأقوال إلى أن المقصود به رفع الحرج في هذه الصورة من جهة أحكام الطلاق التي تتعلق بغير المدخول بها. فإن من لم تمس لا يتعلق طلاقها بما يتعلق بغيرها من اعتبار الحيض والطهر في هذا الباب على الصورة المعروفة في المدخول بها.

وهذا الموضع نافع في تعليم الطالب أن ألفاظ رفع الحرج قد تأتي أحياناً لدفع توهم خاص في صورة مخصوصة، لا لمجرد الإذن المجرد.

\separator

\section{تأويل (وإن طلقتموهن من قبل أن تمسوهن وقد فرضتم لهن فريضة)}
هذه هي الصورة الثانية المفصلة: أن يكون قد فرض لها صداق، ثم وقع الطلاق قبل الدخول. فقرر \rawi{الطبري} أن لها نصف ما فرض، وأن هذا بيان لما أجمل أولاً، وتفريق بين من فرض لها ومن لم يفرض لها.

ونبّه الشيخ هنا إلى فائدة دقيقة: أن القرآن قد يعيد ذكر الصورة لا على وجه التكرار المحض، بل لدفع اللبس، وبيان أن هذه الصورة لها حكم آخر. فليس كل إعادة في التنزيل تكراراً، بل كثير منها من تمام البيان.

\begin{fawaid}[التكرار الظاهر في القرآن قد يكون لتمام البيان لا لحشو اللفظ]
من فقه النظم أن الآية قد تعيد ذكر صورة متقدمة لتفصل فيها حكماً جديداً، أو لتدفع توهماً، أو لتقطع احتمالاً فاسداً. ومن لم يلحظ هذا ظن البيان تكراراً لا ثمرة له.
\end{fawaid}

\section{تأويل (فنصف ما فرضتم إلا أن يعفون)}
بيّن \rawi{الطبري} أن الأصل هنا ثبوت نصف الصداق للمرأة إذا طلقت قبل المسيس بعد تسمية المهر، إلا أن تعفو المرأة عن حقها فتترك ما وجب لها تفضلاً منها، إذا كانت ممن يجوز تصرفها في مالها، أي: بالغة رشيدة.

وهنا يتجلى أدب التفسير والفقه معاً؛ إذ سمي تركها لهذا الحق \textit{عفواً} و\textit{تفضلاً}، لأنه حق ثبت لها، فإسقاطها له ليس واجباً، بل إحسان منها إن فعلته.

\begin{fawaid}[وصف إسقاط الحق بالعفو يربي على معرفة المقامات]
إذا علمت المرأة أو الرجل أن ما يتركه هنا هو حق ثابت له، عرف أن إسقاطه ليس إلزاماً شرعياً، بل مقام فضل وإحسان. وهذا أدق في تربية النفوس من التعبير المجرد بترك المال.
\end{fawaid}

\section{تأويل (أو يعفو الذي بيده عقدة النكاح)}
نقل \rawi{الطبري} في المراد بهذا القائل قولين مشهورين:
\begin{itemize}
    \item فقيل: هو الولي، ولا سيما في حق البكر الصغيرة ونحوها.
    \item وقيل: بل هو الزوج، فيعفو عن حقه ويعطي المرأة الصداق كاملاً.
\end{itemize}

وبيّن الشيخ أن هذا من المواضع التي ينبني فيها اختلاف التفسير على تعيين المخاطب وصاحب الحق في السياق. وهو باب من أبواب الخلاف المتكرر في التفسير: هل يعود الضمير أو الوصف إلى هذا الطرف أو ذاك؟ وما الأثر المترتب على ذلك؟

والذي يظهر من سياق الآية وسياق العفو أن المقصود بالعفو قد يجري على وجهين في النظر: عفو المرأة عن بعض حقها، أو تفضل الزوج بإتمام الصداق. ولذلك كان هذا الباب من أدق أبواب التحرير في التفسير الفقهي.

\begin{fawaid}[تعيين صاحب الضمير أو الوصف قد يغير الحكم الفقهي كله]
من أعظم ما يحتاجه طالب التفسير أن ينتبه إلى أن بعض الخلاف لا يكون في أصل المعنى، بل في تعيين من يعود عليه الوصف أو الضمير. وهذا قد يترتب عليه اختلاف فقهي كامل كما في هذه الآية.
\end{fawaid}

\section{فائدة حديثية: تنبه الطبري لأوهام الرواة}
مرّ في هذا المجلس مثال لطيف من تنبه \rawi{الطبري} إلى خطإ بعض الرواة في الاسم، فصحح الموضع ونبّه عليه. وهذه من الجوانب التي تؤكد أن قراءته ليست مجرد تفسير للمعاني، بل فيها عناية ظاهرة بالأسانيد، وأسماء الرجال، وتمييز ما يقع فيها من وهم.

\begin{fawaid}[العناية بالأسانيد تضبط الفهم كما تضبط الرواية]
من يقرأ الطبري قراءة متأنية يجد أنه لا يكتفي بنقل الآثار، بل يلاحظ ما يقع في الأسماء والأسانيد من تصحيف أو قلب أو وهم. وهذه عناية تفيد في نقد المنقول كما تفيد في صحة الفهم.
\end{fawaid}

\separator

خلاصة هذا اللقاء أن آيات الطلاق قبل المسيس جاءت على غاية من الدقة: فرقت بين من فرض لها ومن لم يفرض لها، وجعلت للمتعة موضعها، ولنصف الصداق موضعه، وفتحت باب الفضل بالعفو من الطرف المستحق. وظهر في ذلك كله جمال صنيع \rawi{الطبري} في ربط اللفظ العربي بالنظم القرآني وبالأثر الفقهي المترتب عليه.
